% ============================================================================
%  2. Arquitectura del sistema
% ============================================================================
\section{Arquitectura del sistema}
\label{sec:arquitectura}

El sistema separa la lógica (\SI{5}{\volt} regulados por el LM7805) de la
potencia (batería de hasta \SI{10}{\volt} hacia el VM del DRV8833), unidas solo
por una masa común. El PIC genera cinco señales hacia el driver: dos PWM de
hardware para la velocidad, dos líneas de dirección y la habilitación nSLEEP
(figura~\ref{fig:arquitectura}).

\begin{figure}[H]
  \centering
  % ============================================================================
%  Diagrama de bloques del sistema (TikZ)
%  Se incluye desde secciones/03-arquitectura.tex con % ============================================================================
%  Diagrama de bloques del sistema (TikZ)
%  Se incluye desde secciones/03-arquitectura.tex con % ============================================================================
%  Diagrama de bloques del sistema (TikZ)
%  Se incluye desde secciones/03-arquitectura.tex con \input{figuras/arquitectura}
% ============================================================================
\begin{tikzpicture}[
    font=\small,
    >={Stealth[length=2.2mm]},
    bloque/.style={draw=gris, rounded corners=3pt, line width=0.6pt,
                   minimum width=3.2cm, minimum height=1.1cm, align=center},
    principal/.style={draw=acento, fill=acentoclaro, rounded corners=3pt,
                      line width=1.2pt, align=left, inner sep=6pt},
    driver/.style={draw=gris, rounded corners=3pt, line width=0.6pt,
                   align=left, inner sep=6pt},
    flecha/.style={->, draw=gris, line width=0.6pt},
    etiqueta/.style={font=\footnotesize, text=gris}
  ]

  % ---------- Microcontrolador ----------
  \node[principal, minimum width=5.2cm, minimum height=9.6cm, anchor=north]
       (pic) at (0,0) {};
  \node[font=\bfseries\large, anchor=north] at ($(pic.north)+(0,-0.25)$) {PIC16F877A};
  \node[anchor=north west, align=left, font=\footnotesize]
       at ($(pic.north west)+(0.2,-1.0)$) {%
       RC2 / CCP1: PWM izquierdo\\[2pt]
       RD0: dirección izquierda\\[2pt]
       RC1 / CCP2: PWM derecho\\[2pt]
       RD1: dirección derecha\\[2pt]
       RD2: nSLEEP};

  % ---------- Columna izquierda ----------
  \node[bloque, anchor=north] (bat) at (-5.6,-0.1)
       {\textbf{Batería}\\ \footnotesize\textcolor{gris}{máx.\ \SI{10}{\volt} DC}};
  \node[bloque, below=0.75cm of bat] (reg)
       {\textbf{LM7805}\\ \footnotesize\textcolor{gris}{\SI{5}{\volt} para la lógica}};
  \node[bloque, below=0.75cm of reg] (xtal)
       {\textbf{Cristal}\\ \footnotesize\textcolor{gris}{\SI{20}{\mega\hertz} (o \SI{4}{\mega\hertz})}};
  \node[bloque, below=0.75cm of xtal] (pk)
       {\textbf{PICkit 3}\\ \footnotesize\textcolor{gris}{programación ICSP}};
  \node[bloque, below=0.75cm of pk] (bt)
       {\textbf{HC-05 / HC-06}\\ \footnotesize\textcolor{gris}{Bluetooth, 9600 baud}};

  % ---------- Columna derecha ----------
  \node[driver, minimum width=4.2cm, minimum height=4.2cm, anchor=north]
       (drv) at (6.4,-0.1) {};
  \node[font=\bfseries\large, anchor=north] at ($(drv.north)+(0,-0.2)$) {DRV8833};
  \node[anchor=north west, align=left, font=\footnotesize]
       at ($(drv.north west)+(0.15,-0.9)$) {%
       AIN1 / AIN2: rueda izq.\\[2pt]
       BIN1 / BIN2: rueda der.\\[2pt]
       nSLEEP: habilitación\\[2pt]
       \textcolor{gris}{VM: 2.7 a \SI{10.8}{\volt}}\\[2pt]
       \textcolor{gris}{\SI{1.5}{\ampere} por canal}};

  \node[bloque, minimum width=4.2cm, below=0.9cm of drv] (mizq)
       {\textbf{Motor izquierdo}\\ \footnotesize\textcolor{gris}{AOUT1 / AOUT2}};
  \node[bloque, minimum width=4.2cm, below=0.75cm of mizq] (mder)
       {\textbf{Motor derecho}\\ \footnotesize\textcolor{gris}{BOUT1 / BOUT2}};

  % ---------- Conexiones ----------
  % Potencia de motores
  \draw[flecha] (bat.north) -- ++(0,0.45) -| node[etiqueta, pos=0.25, above]
       {Vmotor $\leq$ \SI{10}{\volt}} (drv.north);
  \draw[flecha] (bat.south) -- (reg.north);

  % Lógica y periféricos hacia el PIC
  \draw[flecha] (reg.east)  -- node[etiqueta, above] {\SI{5}{\volt}} (reg.east -| pic.west);
  \draw[flecha] (xtal.east) -- (xtal.east -| pic.west);
  \draw[flecha] (pk.east)   -- (pk.east -| pic.west);
  \draw[<->, draw=gris, line width=0.6pt]
       (bt.east) -- node[etiqueta, above] {UART} (bt.east -| pic.west);

  % Pines del PIC alineados con cada flecha que llega
  \node[anchor=west, font=\footnotesize, text=gris] at (xtal.east -| pic.west)
       [xshift=0.2cm, yshift=0.15cm] {OSC1 / OSC2: cristal};
  \node[anchor=west, font=\footnotesize, text=gris] at (xtal.east -| pic.west)
       [xshift=0.2cm, yshift=-0.25cm] {Timer2: base de tiempo PWM};
  \node[anchor=west, font=\footnotesize, text=gris] at (pk.east -| pic.west)
       [xshift=0.2cm] {MCLR, RB6, RB7: ICSP};
  \node[anchor=west, font=\footnotesize, text=gris] at (bt.east -| pic.west)
       [xshift=0.2cm] {RC6 TX / RC7 RX: UART};

  % Señales de control hacia el driver (línea horizontal)
  \coordinate (sal) at ($(drv.west)+(0,0.6)$);
  \draw[flecha] (sal -| pic.east) -- node[etiqueta, above] {5 señales} (sal);

  % Salidas a los motores
  \draw[flecha] (drv.south) -- (mizq.north);
  \draw[flecha] ($(drv.east)+(0,-1.0)$) -- ++(0.4,0) |- (mder.east);
\end{tikzpicture}

% ============================================================================
\begin{tikzpicture}[
    font=\small,
    >={Stealth[length=2.2mm]},
    bloque/.style={draw=gris, rounded corners=3pt, line width=0.6pt,
                   minimum width=3.2cm, minimum height=1.1cm, align=center},
    principal/.style={draw=acento, fill=acentoclaro, rounded corners=3pt,
                      line width=1.2pt, align=left, inner sep=6pt},
    driver/.style={draw=gris, rounded corners=3pt, line width=0.6pt,
                   align=left, inner sep=6pt},
    flecha/.style={->, draw=gris, line width=0.6pt},
    etiqueta/.style={font=\footnotesize, text=gris}
  ]

  % ---------- Microcontrolador ----------
  \node[principal, minimum width=5.2cm, minimum height=9.6cm, anchor=north]
       (pic) at (0,0) {};
  \node[font=\bfseries\large, anchor=north] at ($(pic.north)+(0,-0.25)$) {PIC16F877A};
  \node[anchor=north west, align=left, font=\footnotesize]
       at ($(pic.north west)+(0.2,-1.0)$) {%
       RC2 / CCP1: PWM izquierdo\\[2pt]
       RD0: dirección izquierda\\[2pt]
       RC1 / CCP2: PWM derecho\\[2pt]
       RD1: dirección derecha\\[2pt]
       RD2: nSLEEP};

  % ---------- Columna izquierda ----------
  \node[bloque, anchor=north] (bat) at (-5.6,-0.1)
       {\textbf{Batería}\\ \footnotesize\textcolor{gris}{máx.\ \SI{10}{\volt} DC}};
  \node[bloque, below=0.75cm of bat] (reg)
       {\textbf{LM7805}\\ \footnotesize\textcolor{gris}{\SI{5}{\volt} para la lógica}};
  \node[bloque, below=0.75cm of reg] (xtal)
       {\textbf{Cristal}\\ \footnotesize\textcolor{gris}{\SI{20}{\mega\hertz} (o \SI{4}{\mega\hertz})}};
  \node[bloque, below=0.75cm of xtal] (pk)
       {\textbf{PICkit 3}\\ \footnotesize\textcolor{gris}{programación ICSP}};
  \node[bloque, below=0.75cm of pk] (bt)
       {\textbf{HC-05 / HC-06}\\ \footnotesize\textcolor{gris}{Bluetooth, 9600 baud}};

  % ---------- Columna derecha ----------
  \node[driver, minimum width=4.2cm, minimum height=4.2cm, anchor=north]
       (drv) at (6.4,-0.1) {};
  \node[font=\bfseries\large, anchor=north] at ($(drv.north)+(0,-0.2)$) {DRV8833};
  \node[anchor=north west, align=left, font=\footnotesize]
       at ($(drv.north west)+(0.15,-0.9)$) {%
       AIN1 / AIN2: rueda izq.\\[2pt]
       BIN1 / BIN2: rueda der.\\[2pt]
       nSLEEP: habilitación\\[2pt]
       \textcolor{gris}{VM: 2.7 a \SI{10.8}{\volt}}\\[2pt]
       \textcolor{gris}{\SI{1.5}{\ampere} por canal}};

  \node[bloque, minimum width=4.2cm, below=0.9cm of drv] (mizq)
       {\textbf{Motor izquierdo}\\ \footnotesize\textcolor{gris}{AOUT1 / AOUT2}};
  \node[bloque, minimum width=4.2cm, below=0.75cm of mizq] (mder)
       {\textbf{Motor derecho}\\ \footnotesize\textcolor{gris}{BOUT1 / BOUT2}};

  % ---------- Conexiones ----------
  % Potencia de motores
  \draw[flecha] (bat.north) -- ++(0,0.45) -| node[etiqueta, pos=0.25, above]
       {Vmotor $\leq$ \SI{10}{\volt}} (drv.north);
  \draw[flecha] (bat.south) -- (reg.north);

  % Lógica y periféricos hacia el PIC
  \draw[flecha] (reg.east)  -- node[etiqueta, above] {\SI{5}{\volt}} (reg.east -| pic.west);
  \draw[flecha] (xtal.east) -- (xtal.east -| pic.west);
  \draw[flecha] (pk.east)   -- (pk.east -| pic.west);
  \draw[<->, draw=gris, line width=0.6pt]
       (bt.east) -- node[etiqueta, above] {UART} (bt.east -| pic.west);

  % Pines del PIC alineados con cada flecha que llega
  \node[anchor=west, font=\footnotesize, text=gris] at (xtal.east -| pic.west)
       [xshift=0.2cm, yshift=0.15cm] {OSC1 / OSC2: cristal};
  \node[anchor=west, font=\footnotesize, text=gris] at (xtal.east -| pic.west)
       [xshift=0.2cm, yshift=-0.25cm] {Timer2: base de tiempo PWM};
  \node[anchor=west, font=\footnotesize, text=gris] at (pk.east -| pic.west)
       [xshift=0.2cm] {MCLR, RB6, RB7: ICSP};
  \node[anchor=west, font=\footnotesize, text=gris] at (bt.east -| pic.west)
       [xshift=0.2cm] {RC6 TX / RC7 RX: UART};

  % Señales de control hacia el driver (línea horizontal)
  \coordinate (sal) at ($(drv.west)+(0,0.6)$);
  \draw[flecha] (sal -| pic.east) -- node[etiqueta, above] {5 señales} (sal);

  % Salidas a los motores
  \draw[flecha] (drv.south) -- (mizq.north);
  \draw[flecha] ($(drv.east)+(0,-1.0)$) -- ++(0.4,0) |- (mder.east);
\end{tikzpicture}

% ============================================================================
\begin{tikzpicture}[
    font=\small,
    >={Stealth[length=2.2mm]},
    bloque/.style={draw=gris, rounded corners=3pt, line width=0.6pt,
                   minimum width=3.2cm, minimum height=1.1cm, align=center},
    principal/.style={draw=acento, fill=acentoclaro, rounded corners=3pt,
                      line width=1.2pt, align=left, inner sep=6pt},
    driver/.style={draw=gris, rounded corners=3pt, line width=0.6pt,
                   align=left, inner sep=6pt},
    flecha/.style={->, draw=gris, line width=0.6pt},
    etiqueta/.style={font=\footnotesize, text=gris}
  ]

  % ---------- Microcontrolador ----------
  \node[principal, minimum width=5.2cm, minimum height=9.6cm, anchor=north]
       (pic) at (0,0) {};
  \node[font=\bfseries\large, anchor=north] at ($(pic.north)+(0,-0.25)$) {PIC16F877A};
  \node[anchor=north west, align=left, font=\footnotesize]
       at ($(pic.north west)+(0.2,-1.0)$) {%
       RC2 / CCP1: PWM izquierdo\\[2pt]
       RD0: dirección izquierda\\[2pt]
       RC1 / CCP2: PWM derecho\\[2pt]
       RD1: dirección derecha\\[2pt]
       RD2: nSLEEP};

  % ---------- Columna izquierda ----------
  \node[bloque, anchor=north] (bat) at (-5.6,-0.1)
       {\textbf{Batería}\\ \footnotesize\textcolor{gris}{máx.\ \SI{10}{\volt} DC}};
  \node[bloque, below=0.75cm of bat] (reg)
       {\textbf{LM7805}\\ \footnotesize\textcolor{gris}{\SI{5}{\volt} para la lógica}};
  \node[bloque, below=0.75cm of reg] (xtal)
       {\textbf{Cristal}\\ \footnotesize\textcolor{gris}{\SI{20}{\mega\hertz} (o \SI{4}{\mega\hertz})}};
  \node[bloque, below=0.75cm of xtal] (pk)
       {\textbf{PICkit 3}\\ \footnotesize\textcolor{gris}{programación ICSP}};
  \node[bloque, below=0.75cm of pk] (bt)
       {\textbf{HC-05 / HC-06}\\ \footnotesize\textcolor{gris}{Bluetooth, 9600 baud}};

  % ---------- Columna derecha ----------
  \node[driver, minimum width=4.2cm, minimum height=4.2cm, anchor=north]
       (drv) at (6.4,-0.1) {};
  \node[font=\bfseries\large, anchor=north] at ($(drv.north)+(0,-0.2)$) {DRV8833};
  \node[anchor=north west, align=left, font=\footnotesize]
       at ($(drv.north west)+(0.15,-0.9)$) {%
       AIN1 / AIN2: rueda izq.\\[2pt]
       BIN1 / BIN2: rueda der.\\[2pt]
       nSLEEP: habilitación\\[2pt]
       \textcolor{gris}{VM: 2.7 a \SI{10.8}{\volt}}\\[2pt]
       \textcolor{gris}{\SI{1.5}{\ampere} por canal}};

  \node[bloque, minimum width=4.2cm, below=0.9cm of drv] (mizq)
       {\textbf{Motor izquierdo}\\ \footnotesize\textcolor{gris}{AOUT1 / AOUT2}};
  \node[bloque, minimum width=4.2cm, below=0.75cm of mizq] (mder)
       {\textbf{Motor derecho}\\ \footnotesize\textcolor{gris}{BOUT1 / BOUT2}};

  % ---------- Conexiones ----------
  % Potencia de motores
  \draw[flecha] (bat.north) -- ++(0,0.45) -| node[etiqueta, pos=0.25, above]
       {Vmotor $\leq$ \SI{10}{\volt}} (drv.north);
  \draw[flecha] (bat.south) -- (reg.north);

  % Lógica y periféricos hacia el PIC
  \draw[flecha] (reg.east)  -- node[etiqueta, above] {\SI{5}{\volt}} (reg.east -| pic.west);
  \draw[flecha] (xtal.east) -- (xtal.east -| pic.west);
  \draw[flecha] (pk.east)   -- (pk.east -| pic.west);
  \draw[<->, draw=gris, line width=0.6pt]
       (bt.east) -- node[etiqueta, above] {UART} (bt.east -| pic.west);

  % Pines del PIC alineados con cada flecha que llega
  \node[anchor=west, font=\footnotesize, text=gris] at (xtal.east -| pic.west)
       [xshift=0.2cm, yshift=0.15cm] {OSC1 / OSC2: cristal};
  \node[anchor=west, font=\footnotesize, text=gris] at (xtal.east -| pic.west)
       [xshift=0.2cm, yshift=-0.25cm] {Timer2: base de tiempo PWM};
  \node[anchor=west, font=\footnotesize, text=gris] at (pk.east -| pic.west)
       [xshift=0.2cm] {MCLR, RB6, RB7: ICSP};
  \node[anchor=west, font=\footnotesize, text=gris] at (bt.east -| pic.west)
       [xshift=0.2cm] {RC6 TX / RC7 RX: UART};

  % Señales de control hacia el driver (línea horizontal)
  \coordinate (sal) at ($(drv.west)+(0,0.6)$);
  \draw[flecha] (sal -| pic.east) -- node[etiqueta, above] {5 señales} (sal);

  % Salidas a los motores
  \draw[flecha] (drv.south) -- (mizq.north);
  \draw[flecha] ($(drv.east)+(0,-1.0)$) -- ++(0.4,0) |- (mder.east);
\end{tikzpicture}

  \caption{Arquitectura del sistema: control, potencia, actuadores y
           comunicaciones.}
  \label{fig:arquitectura}
\end{figure}

\subsection{Asignación de pines}

\begin{table}[H]
  \centering
  \caption{Asignación de pines del PIC16F877A y justificación.}
  \label{tab:pines}
  \small
  \begin{tabularx}{\textwidth}{@{}l l l L@{}}
    \toprule
    \textbf{Pin PIC} & \textbf{Función} & \textbf{Conectado a} & \textbf{Justificación} \\
    \midrule
    RC2 / CCP1 (17) & PWM motor izq.   & AIN1 & Único pin de salida del módulo CCP1; PWM por hardware sin carga de CPU. \\
    RC1 / CCP2 (16) & PWM motor der.   & BIN1 & Único pin de salida de CCP2; comparte Timer2 con CCP1, así ambas ruedas tienen la misma frecuencia y fase. \\
    RD0 (19)        & Dirección izq.   & AIN2 & Puerto D libre de funciones especiales que interfieran. \\
    RD1 (20)        & Dirección der.   & BIN2 & Ídem. \\
    RD2 (21)        & Habilitación     & nSLEEP & Permite dormir el driver hasta terminar la configuración y aplicar el pulso de activación. \\
    RC6 / TX (25)   & UART TX          & RXD HC-05 & Pin fijo del USART; un divisor 1k/2k baja \SI{5}{\volt} a $\approx$\SI{3.3}{\volt}. \\
    RC7 / RX (26)   & UART RX          & TXD HC-05 & Pin fijo del USART con interrupción de recepción. \\
    RB6, RB7, MCLR  & ICSP             & PICkit 3 & Reservados para programar en circuito sin retirar el PIC. \\
    \bottomrule
  \end{tabularx}
\end{table}

\subsection{Componentes principales}

PIC16F877A, módulo DRV8833, dos motorreductores DC (caja amarilla tipo TT),
módulo Bluetooth HC-05/HC-06, regulador LM7805, cristal de cuarzo, batería de
máximo \SI{10}{\volt} DC y programador PICkit~3.
